\IfFileExists{stacks-project.cls}{%
\documentclass{stacks-project}
}{%
\documentclass{amsart}
}

% For dealing with references we use the comment environment
\usepackage{verbatim}
\newenvironment{reference}{\comment}{\endcomment}
%\newenvironment{reference}{}{}
\newenvironment{slogan}{\comment}{\endcomment}
\newenvironment{history}{\comment}{\endcomment}

% For commutative diagrams we use Xy-pic
\usepackage[all]{xy}

% We use 2cell for 2-commutative diagrams.
\xyoption{2cell}
\UseAllTwocells

% We use multicol for the list of chapters between chapters
\usepackage{multicol}

% This is generally recommended for better output
\usepackage{lmodern}
\usepackage[T1]{fontenc}

% For cross-file-references
\usepackage{xr-hyper}

% Package for hypertext links:
\usepackage{hyperref}

% For any local file, say "hello.tex" you want to link to please
% use \externaldocument[hello-]{hello}
\externaldocument[introduction-]{introduction}
\externaldocument[conventions-]{conventions}
\externaldocument[sets-]{sets}
\externaldocument[categories-]{categories}
\externaldocument[topology-]{topology}
\externaldocument[sheaves-]{sheaves}
\externaldocument[sites-]{sites}
\externaldocument[stacks-]{stacks}
\externaldocument[fields-]{fields}
\externaldocument[algebra-]{algebra}
\externaldocument[brauer-]{brauer}
\externaldocument[homology-]{homology}
\externaldocument[derived-]{derived}
\externaldocument[simplicial-]{simplicial}
\externaldocument[more-algebra-]{more-algebra}
\externaldocument[smoothing-]{smoothing}
\externaldocument[modules-]{modules}
\externaldocument[sites-modules-]{sites-modules}
\externaldocument[injectives-]{injectives}
\externaldocument[cohomology-]{cohomology}
\externaldocument[sites-cohomology-]{sites-cohomology}
\externaldocument[dga-]{dga}
\externaldocument[dpa-]{dpa}
\externaldocument[sdga-]{sdga}
\externaldocument[hypercovering-]{hypercovering}
\externaldocument[schemes-]{schemes}
\externaldocument[constructions-]{constructions}
\externaldocument[properties-]{properties}
\externaldocument[morphisms-]{morphisms}
\externaldocument[coherent-]{coherent}
\externaldocument[divisors-]{divisors}
\externaldocument[limits-]{limits}
\externaldocument[varieties-]{varieties}
\externaldocument[topologies-]{topologies}
\externaldocument[descent-]{descent}
\externaldocument[perfect-]{perfect}
\externaldocument[more-morphisms-]{more-morphisms}
\externaldocument[flat-]{flat}
\externaldocument[groupoids-]{groupoids}
\externaldocument[more-groupoids-]{more-groupoids}
\externaldocument[etale-]{etale}
\externaldocument[chow-]{chow}
\externaldocument[intersection-]{intersection}
\externaldocument[pic-]{pic}
\externaldocument[weil-]{weil}
\externaldocument[adequate-]{adequate}
\externaldocument[dualizing-]{dualizing}
\externaldocument[duality-]{duality}
\externaldocument[discriminant-]{discriminant}
\externaldocument[derham-]{derham}
\externaldocument[local-cohomology-]{local-cohomology}
\externaldocument[algebraization-]{algebraization}
\externaldocument[curves-]{curves}
\externaldocument[resolve-]{resolve}
\externaldocument[models-]{models}
\externaldocument[functors-]{functors}
\externaldocument[equiv-]{equiv}
\externaldocument[pione-]{pione}
\externaldocument[etale-cohomology-]{etale-cohomology}
\externaldocument[proetale-]{proetale}
\externaldocument[relative-cycles-]{relative-cycles}
\externaldocument[more-etale-]{more-etale}
\externaldocument[trace-]{trace}
\externaldocument[crystalline-]{crystalline}
\externaldocument[spaces-]{spaces}
\externaldocument[spaces-properties-]{spaces-properties}
\externaldocument[spaces-morphisms-]{spaces-morphisms}
\externaldocument[decent-spaces-]{decent-spaces}
\externaldocument[spaces-cohomology-]{spaces-cohomology}
\externaldocument[spaces-limits-]{spaces-limits}
\externaldocument[spaces-divisors-]{spaces-divisors}
\externaldocument[spaces-over-fields-]{spaces-over-fields}
\externaldocument[spaces-topologies-]{spaces-topologies}
\externaldocument[spaces-descent-]{spaces-descent}
\externaldocument[spaces-perfect-]{spaces-perfect}
\externaldocument[spaces-more-morphisms-]{spaces-more-morphisms}
\externaldocument[spaces-flat-]{spaces-flat}
\externaldocument[spaces-groupoids-]{spaces-groupoids}
\externaldocument[spaces-more-groupoids-]{spaces-more-groupoids}
\externaldocument[bootstrap-]{bootstrap}
\externaldocument[spaces-pushouts-]{spaces-pushouts}
\externaldocument[spaces-chow-]{spaces-chow}
\externaldocument[groupoids-quotients-]{groupoids-quotients}
\externaldocument[spaces-more-cohomology-]{spaces-more-cohomology}
\externaldocument[spaces-simplicial-]{spaces-simplicial}
\externaldocument[spaces-duality-]{spaces-duality}
\externaldocument[formal-spaces-]{formal-spaces}
\externaldocument[restricted-]{restricted}
\externaldocument[spaces-resolve-]{spaces-resolve}
\externaldocument[formal-defos-]{formal-defos}
\externaldocument[defos-]{defos}
\externaldocument[cotangent-]{cotangent}
\externaldocument[examples-defos-]{examples-defos}
\externaldocument[algebraic-]{algebraic}
\externaldocument[examples-stacks-]{examples-stacks}
\externaldocument[stacks-sheaves-]{stacks-sheaves}
\externaldocument[criteria-]{criteria}
\externaldocument[artin-]{artin}
\externaldocument[quot-]{quot}
\externaldocument[stacks-properties-]{stacks-properties}
\externaldocument[stacks-morphisms-]{stacks-morphisms}
\externaldocument[stacks-limits-]{stacks-limits}
\externaldocument[stacks-cohomology-]{stacks-cohomology}
\externaldocument[stacks-perfect-]{stacks-perfect}
\externaldocument[stacks-introduction-]{stacks-introduction}
\externaldocument[stacks-more-morphisms-]{stacks-more-morphisms}
\externaldocument[stacks-geometry-]{stacks-geometry}
\externaldocument[moduli-]{moduli}
\externaldocument[moduli-curves-]{moduli-curves}
\externaldocument[examples-]{examples}
\externaldocument[exercises-]{exercises}
\externaldocument[guide-]{guide}
\externaldocument[desirables-]{desirables}
\externaldocument[coding-]{coding}
\externaldocument[obsolete-]{obsolete}
\externaldocument[fdl-]{fdl}
\externaldocument[index-]{index}

% Theorem environments.
%
\theoremstyle{plain}
\newtheorem{theorem}[subsection]{Theorem}
\newtheorem{proposition}[subsection]{Proposition}
\newtheorem{lemma}[subsection]{Lemma}

\theoremstyle{definition}
\newtheorem{definition}[subsection]{Definition}
\newtheorem{example}[subsection]{Example}
\newtheorem{exercise}[subsection]{Exercise}
\newtheorem{situation}[subsection]{Situation}

\theoremstyle{remark}
\newtheorem{remark}[subsection]{Remark}
\newtheorem{remarks}[subsection]{Remarks}

\numberwithin{equation}{subsection}

% Macros
%
\def\lim{\mathop{\mathrm{lim}}\nolimits}
\def\colim{\mathop{\mathrm{colim}}\nolimits}
\def\Spec{\mathop{\mathrm{Spec}}}
\def\Hom{\mathop{\mathrm{Hom}}\nolimits}
\def\Ext{\mathop{\mathrm{Ext}}\nolimits}
\def\SheafHom{\mathop{\mathcal{H}\!\mathit{om}}\nolimits}
\def\SheafExt{\mathop{\mathcal{E}\!\mathit{xt}}\nolimits}
\def\Sch{\mathit{Sch}}
\def\Mor{\mathop{\mathrm{Mor}}\nolimits}
\def\Ob{\mathop{\mathrm{Ob}}\nolimits}
\def\Sh{\mathop{\mathit{Sh}}\nolimits}
\def\NL{\mathop{N\!L}\nolimits}
\def\CH{\mathop{\mathrm{CH}}\nolimits}
\def\proetale{{pro\text{-}\acute{e}tale}}
\def\etale{{\acute{e}tale}}
\def\QCoh{\mathit{QCoh}}
\def\Ker{\mathop{\mathrm{Ker}}}
\def\Im{\mathop{\mathrm{Im}}}
\def\Coker{\mathop{\mathrm{Coker}}}
\def\Coim{\mathop{\mathrm{Coim}}}

% Boxtimes
%
\DeclareMathSymbol{\boxtimes}{\mathbin}{AMSa}{"02}

%
% Macros for moduli stacks/spaces
%
\def\QCohstack{\mathcal{QC}\!\mathit{oh}}
\def\Cohstack{\mathcal{C}\!\mathit{oh}}
\def\Spacesstack{\mathcal{S}\!\mathit{paces}}
\def\Quotfunctor{\mathrm{Quot}}
\def\Hilbfunctor{\mathrm{Hilb}}
\def\Curvesstack{\mathcal{C}\!\mathit{urves}}
\def\Polarizedstack{\mathcal{P}\!\mathit{olarized}}
\def\Complexesstack{\mathcal{C}\!\mathit{omplexes}}
% \Pic is the operator that assigns to X its picard group, usage \Pic(X)
% \Picardstack_{X/B} denotes the Picard stack of X over B
% \Picardfunctor_{X/B} denotes the Picard functor of X over B
\def\Pic{\mathop{\mathrm{Pic}}\nolimits}
\def\Picardstack{\mathcal{P}\!\mathit{ic}}
\def\Picardfunctor{\mathrm{Pic}}
\def\Deformationcategory{\mathcal{D}\!\mathit{ef}}


% OK, start here.
%
\begin{document}

\title{Topology}


\maketitle

\phantomsection
\label{section-phantom}

\tableofcontents

\section{Introduction}
\label{section-introduction}

\noindent
Basic topology will be explained in this document.
A reference is \cite{Engelking}.

\section{Basic notions}
\label{section-topology-basic}

\noindent
The following is a list of basic notions in topology. Some of these notions
are discussed in more detail in the text that follows and some
are defined in the list, but others are considered basic and will not
be defined. If you are not familiar with most of the italicized concepts,
then we suggest looking at an introductory text on topology before
continuing.

\begin{enumerate}
\item
\label{item-space}
$X$ is a {\it topological space},
\item
\label{item-point}
$x\in X$ is a {\it point},
\item
\label{item-locally-closed}
$E \subset X$ is a {\it locally closed} subset,
\item
\label{item-closed-point}
$x\in X$ is a {\it closed point},
\item
\label{item-dense}
$E \subset X$ is a {\it dense} subset,
\item
\label{item-continuous}
$f : X_1 \to X_2$ is {\it continuous},
\item
\label{item-upper-semi-continuous}
an extended real function $f : X \to \mathbf{R} \cup \{\infty, -\infty\}$
is {\it upper semi-continuous} if $\{x \in X \mid f(x) < a\}$ is open for
all $a \in \mathbf{R}$,
\item
\label{item-lower-semi-continuous}
an extended real function $f : X \to \mathbf{R} \cup \{\infty, -\infty\}$
is {\it lower semi-continuous} if $\{x \in X \mid f(x) > a\}$ is open for
all $a \in \mathbf{R}$,
\item a continuous map of spaces $f : X \to Y$ is
{\it open} if $f(U)$ is open in $Y$ for $U \subset X$ open,
\item a continuous map of spaces $f : X \to Y$ is
{\it closed} if $f(Z)$ is closed in $Y$ for $Z \subset X$ closed,
\item
\label{item-neighbourhood}
a {\it neighbourhood of $x \in X$} is any subset
$E \subset X$ which contains an open subset that
contains $x$,
\item
\label{item-induced-topology}
the {\it induced topology} on a subset $E \subset X$,
\item
\label{item-covering}
$\mathcal{U} : U = \bigcup_{i \in I} U_i$ is an
{\it open covering of} $U$ (note: we allow any $U_i$ to be empty
and we even allow, in case $U$ is empty, the empty set for $I$),
\item
\label{item-subcover}
a {\it subcover} of a covering as in (\ref{item-covering})
is an open covering $\mathcal{U}' : U = \bigcup_{i \in I'} U_i$
where $I' \subset I$,
\item
\label{item-refinement}
the open covering $\mathcal{V}$ is a {\it refinement}
of the open covering $\mathcal{U}$ (if
$\mathcal{V} : U = \bigcup_{j \in J} V_j$ and
$\mathcal{U} : U = \bigcup_{i \in I} U_i$
this means each $V_j$ is completely contained in one of the $U_i$),
\item
\label{item-fundamental-system}
{\it $\{ E_i \}_{i \in I}$ is a fundamental system of neighbourhoods
of $x$ in $X$},
\item
\label{item-Hausdorff}
a topological space $X$ is called {\it Hausdorff} or {\it separated}
if and only if for every distinct pair of points $x, y \in X$ there exist
disjoint opens $U, V \subset X$ such that $x \in U$, $y \in V$,
\item the {\it product} of two topological spaces,
\label{item-product}
\item
\label{item-fibre-product}
the {\it fibre product $X \times_Y Z$} of a pair of continuous maps
$f : X \to Y$ and $g : Z \to Y$,
\item
\label{item-discrete-indiscrete}
the {\it discrete topology} and the {\it indiscrete topology} on a set,
\item etc.
\end{enumerate}



\section{Hausdorff spaces}
\label{section-Hausdorff}

\noindent
The category of topological spaces has finite products.

\begin{lemma}
\label{lemma-Hausdorff}
Let $X$ be a topological space. The following are equivalent:
\begin{enumerate}
\item $X$ is Hausdorff,
\item the diagonal $\Delta(X) \subset X \times X$ is closed.
\end{enumerate}
\end{lemma}

\begin{proof}
We suppose that $X$ is Hausdorff. Let $(x, y) \not \in \Delta(X)$, i.e.,
$x \neq y$. There are $U$ and $V$ disjoint open sets of $X$
such that $x \in U$ and $y \in V$. This implies that $U \times V \subset
(X \times X) \setminus \Delta(X) $. This shows that $ (X \times X) \setminus
\Delta(X)$ is an open set of $ X \times X$ which is equivalent to saying that
the diagonal $\Delta(X) \subset X \times X$ is closed in $X \times X$.
The converse is similar:
The complement $(X \times X) \setminus \Delta(X)$ consists
precisely of $(x, y) \in X\times X$ with $ x \neq y$. Thus, if $ \Delta(X)
\subset X \times X$ is closed, then, by the definition of the product topology,
for every such $ (x, y)$, there are opens $U, V \subset X$ with $(x, y) \in
U \times V$ and $(U \times V)\cap \Delta(X) = \emptyset$. In other words, with
$x \in U$ and $y \in V$ such that $U \cap V = \emptyset$.
\end{proof}

\begin{lemma}
\label{lemma-graph-closed}
\begin{slogan}
Graphs of maps to Hausdorff spaces are closed.
\end{slogan}
Let $f : X \to Y$ be a continuous map of topological spaces.
If $Y$ is Hausdorff, then the graph of $f$ is closed in $X \times Y$.
\end{lemma}

\begin{proof}
The graph is the inverse image of the diagonal under the map
$X \times Y \to Y \times Y$. Thus the lemma follows from
Lemma \ref{lemma-Hausdorff}.
\end{proof}

\begin{lemma}
\label{lemma-section-closed}
Let $f : X \to Y$ be a continuous map of topological spaces.
Let $s : Y \to X$ be a continuous map such that $f \circ s = \text{id}_Y$.
If $X$ is Hausdorff, then $s(Y)$ is closed.
\end{lemma}

\begin{proof}
This follows from Lemma \ref{lemma-Hausdorff} as
$s(Y) = \{x \in X \mid x = s(f(x))\}$.
\end{proof}

\begin{lemma}
\label{lemma-fibre-product-closed}
Let $X \to Z$ and $Y \to Z$ be continuous maps of topological spaces.
If $Z$ is Hausdorff, then $X \times_Z Y$ is closed in $X \times Y$.
\end{lemma}

\begin{proof}
This follows from Lemma \ref{lemma-Hausdorff} as
$X \times_Z Y$ is the inverse image of $\Delta(Z)$
under $X \times Y \to Z \times Z$.
\end{proof}




\section{Separated maps}
\label{section-separated}

\end{document}
